%% ma: L12-B17-0022
%% lop: 12
%% chuong: 5
%% bai: 17
%% dang: 12.17.14
%% loai: DS
%% muc: ['NB', 'TH', 'TH', 'VD']
%% dap_an: "SĐSĐ"
%% nguon: "(NBV)-10 ĐỀ KIỂM TRA HK2-TOÁN 12-DUNG CHUNG 3 SÁCH.pdf (H04, Đề số 4), Phần 2 Câu 1"
%% kien_thuc: "Bài 7 (tọa độ trung điểm, trọng tâm, khoảng cách), Bài 17 (phương trình mặt cầu)"
%% trang_thai: cho-duyet
%% ly_do: "Dạng mới đề xuất 12.17.14 (Phương trình mặt cầu theo tâm, điểm đi qua, đường kính, trọng tâm tam giác (Đúng/Sai)), chờ giáo viên duyệt."
%% ghi_chu: "Đáp án gốc SĐSĐ khớp. Ý c) đề in $(z-3)^2$, đúng phải là (z-3/2)^2 nên ý c) sai."
\begin{ex}
Trong không gian $Oxyz$, cho ba điểm $A(1;0;0)$, $B(0;2;0)$ và $C(0;0;3)$.
\choiceTF
{Mặt cầu tâm $B$, bán kính $R=3$ có phương trình là $x^2+(y-2)^2+z^2=3$.}
{Mặt cầu tâm $A$, đi qua $B$ có phương trình là $(x-1)^2+y^2+z^2=5$.}
{Mặt cầu nhận $BC$ làm đường kính có phương trình là $x^2+(y-1)^2+(z-3)^2=\dfrac{13}{4}$.}
{Mặt cầu tâm $O$ và có bán kính $OG$ với $G$ là trọng tâm tam giác $ABC$ có phương trình là $x^2+y^2+z^2=\dfrac{14}{9}$.}
\loigiai{a) Sai: phương trình đúng là $x^2+(y-2)^2+z^2=9$.

b) Đúng: $R=AB=\sqrt{1+4}=\sqrt5$ nên $(x-1)^2+y^2+z^2=5$.

c) Sai: tâm là trung điểm $I\left(0;1;\dfrac32\right)$ của $BC$, $R=\dfrac{BC}2=\dfrac{\sqrt{13}}2$ nên phương trình là $x^2+(y-1)^2+\left(z-\dfrac32\right)^2=\dfrac{13}4$.

d) Đúng: $G\left(\dfrac13;\dfrac23;1\right)$, $OG^2=\dfrac19+\dfrac49+1=\dfrac{14}9$.}
\end{ex}
